% from ega4-2-fr.tex
\oldpage[IV-2]{107}

Cela résulte aussitôt des relations (5.7.9.1) et (5.7.9.2) et de la
définition de la propriété $(S_k)$.

\begin{env}[5.7.12]
\label{IV.2.s5.7.12-fr}
Conformément aux définitions de (5.7.1), étant donnés un anneau noethérien
quelconque $A$ et un $A$-module $M$ de type fini, on définit
$\operatorname{coprof}_A(M)$ comme égal à
\[
\operatorname{coprof}(M)=\sup_{x\in X}
\operatorname{coprof}_{A_x}(M_x),\qquad X=\operatorname{Spec}(A);
\]
nous verrons plus loin (6.11.5) que cette définition coïncide avec celle de
(0, 16.4.9) lorsque $A$ est un anneau noethérien local.
\end{env}

\begin{corollary}[5.7.13]
\label{IV.2.s5.7.13-fr}
Soient $A$, $B$ deux anneaux noethériens, $\rho:A\to B$ un homomorphisme
d'anneaux, $M$ un $B$-module tel que $M_{[\rho]}$ soit un $A$-module de type
fini. Alors on a
\begin{equation}
\operatorname{coprof}_A(M_{[\rho]})\geq\operatorname{coprof}_B(M).
\tag{5.7.13.1}\label{IV.2.s5.7.13.1-fr}
\end{equation}
\end{corollary}

Cela résulte de la définition précédente et des relations (5.7.9.1) et
(5.7.9.2).

\subsection{\texorpdfstring{Préschémas réguliers et propriété $(R_k)$. Critère de normalité de Serre}{Préschémas réguliers et propriété (R\string_k). Critère de normalité de Serre}}

\begin{env}[5.8.1]
\label{IV.2.s5.8.1-fr}
Rappelons $(0_{\mathrm I},4.1.4)$ qu'un espace annelé
$(X,\mathcal O_X)$ est dit \emph{régulier} en un point $x\in X$, si
$\mathcal O_x$ est un anneau régulier. Quand il s'agira de préschémas, nous
n'utiliserons cette terminologie dans ce chapitre que lorsque $X$ est
localement noethérien.
\end{env}

\begin{definition}[5.8.2]
\label{IV.2.s5.8.2-fr}
On dit qu'un préschéma localement noethérien $X$ est régulier en codimension
$\leq k$, ou possède la propriété $(R_k)$ si l'ensemble des points où $X$
n'est pas régulier est de codimension $>k$ (autrement dit (5.1.3), si, pour
tout $x\in X$, la relation $\dim(\mathcal O_x)\leq k$ entraîne que
$\mathcal O_x$ est régulier).
\end{definition}

Dire que $X$ est régulier signifie que $X$ possède la propriété $(R_k)$ pour
tout $k$.

Si $X=\operatorname{Spec}(A)$, où $A$ est un anneau noethérien, on dira que
$A$ possède la propriété $(R_k)$ si $X$ possède cette propriété; dire que
$X$ est régulier signifie que l'anneau $A$ est régulier (0, 17.3.6). Pour un
préschéma localement noethérien quelconque $X$, on dira que $X$ possède la
propriété $(R_k)$ en un point $x\in X$ si l'anneau local $\mathcal O_x$
possède la propriété $(R_k)$; cela veut donc dire que pour toute générisation
$x'$ de $x$ dans $X$, la relation $\dim(\mathcal O_{x'})\leq k$ entraîne que
$\mathcal O_{x'}$ est un anneau local régulier. Dire que $X$ est régulier en
un point $x$ équivaut à dire que $X$ vérifie la propriété $(R_n)$ pour tout
$n\geq0$ au point $x$, en vertu de (0, 17.3.6).

\begin{proposition}[5.8.3]
\label{IV.2.s5.8.3-fr}
Si $k$ est un corps, $X$ un préschéma localement de type fini sur $k$, alors,
pour tout $x\in X$, il existe un voisinage ouvert de $x$ dans $X$ isomorphe à
un sous-schéma d'un $k$-schéma régulier.
\end{proposition}

En effet, il y a un voisinage ouvert affine $U$ de $x$ isomorphe à un
$k$-schéma de la forme $\operatorname{Spec}(A)$, où $A$ est une $k$-algèbre
de type fini; $A$ est par suite isomorphe à un quotient d'une algèbre de
polynômes $k[T_1,\ldots,T_n]$, et on sait que cette dernière est un anneau
régulier (0, 17.3.7).

\begin{env}[5.8.4]
\label{IV.2.s5.8.4-fr}
En vertu de (5.8.2), dire que $X$ possède la propriété $(R_0)$ signifie que
pour tout point maximal $x$ de $X$, l'anneau $\mathcal O_x$ est un corps
(0, 17.1.4), autrement dit que $X$ est réduit en ce point. Comme l'ensemble
$U$ des $x\in X$ où $X$ est réduit est ouvert (le

